\section{形变理论：Taylor--Wiles 方法的技术核心}
\label{sec:ch10-deformation}
% ============================================================================

\providecommand{\rhobar}{\bar{\rho}}

如果说上一节回答的是“Wiles 想证明什么”，
那么这一节回答的就是“他究竟拿什么来证明”。
形变理论之所以成为整个证明的技术心脏，
是因为它把一个本来非常刚性的 mod $\ell$ 表示问题
转换成了一个可以做维数计数、切空间比较与极限构造的局部代数问题。
这一步看似抽象，
实际上极富直觉：
我们先把一个 residual 表示钉住，
然后系统地问——它能以多少种方式被提升到 $\ell$-进世界里？

\textbf{我们遇到了什么问题？}
假设我们已经知道某个 residual 表示
\[
  \rhobar\colon G_\Q\to \GL_2(\F_\ell)
\]
是模的。
接下来要证明某一条具体的 $\ell$-进提升 $\rho$ 也是模的。
若只盯着这条单独的提升，
几乎无从下手；
因为它混杂着在每个素数处的局部条件，
还带着行列式、半稳定性、有限平坦性等结构。
Wiles 的关键观察是：
与其直接抓住这一个点，
不如把所有满足同样局部条件的提升一起收集成一个“空间”，
再把这个空间与模形式一侧的 Hecke 空间比较。

\textbf{核心思想是什么？}
把完备 Noether 局部 $\Lambda$-代数范畴记为 $\mathrm{CNL}_\Lambda$。
对每个对象 $A$，
考虑所有满足规定局部条件的提升
\[
  \rho_A\colon G_\Q\to \GL_2(A)
\]
的严格同构类。
这给出一个形变函子
\[
  \mathcal F\colon \mathrm{CNL}_\Lambda\to \mathrm{Sets}.
\]
若它可表，则存在一个完备局部环 $R$ 使得
\[
  \mathcal F(A)\cong \Hom_{\mathrm{CNL}_\Lambda}(R,A).
\]
也就是说，$R$ 真正把“全部允许提升”参数化了。
模形式一侧的 Hecke 环 $T$ 则参数化“全部允许的模提升”。
因此 $R=T$ 的意义非常直白：
允许的提升与模的提升一样多，而且一一对应。

\textbf{引入之后能得到什么？}
一旦把问题压缩到局部环同构，
证明策略就从“构造模形式”转变为“比较两个局部环的大小”。
环的大小可以通过切空间、关系数、长度、Fitting ideal 与 congruence ideal 来测量。
Selmer 群控制 $R$ 的切空间，
模形式之间的同余控制 $T$ 的 congruence ideal；
Taylor--Wiles 再用辅助素数把这两边同时扩展到一个更容易比较的无限层。
这就是为什么 Wiles 证明虽然写在 Galois 表示语言里，
却处处伴随着交换代数与同调代数。

\input{figures/ch10-deformation-space}

先把基本对象写清楚。

\begin{definition}[形变问题]
设 $\rhobar\colon G_\Q\to \GL_2(k)$，其中 $k$ 是有限域。
若 $A$ 是带剩余域 $k$ 的完备 Noether 局部环，
则一个 \textbf{$A$-值提升} 是连续表示
\[
  \rho\colon G_\Q\to \GL_2(A)
\]
使得模去极大理想后恢复 $\rhobar$。
两个提升若相差一个同余于单位阵的共轭，称为严格同构。
给定行列式与局部条件后，
所有这样的严格同构类构成形变函子 $\mathcal F(A)$。
\end{definition}

\begin{theorem}[Mazur 的可表性定理]
\label{thm:ch10-mazur}
若 $\rhobar$ 满足通常的不可约性与奇性条件，
则带有给定局部条件的形变函子 $\mathcal F$ 可由某个完备 Noether 局部环 $R$ 表示。
等价地，存在一个泛形变
\[
  \rho^{\mathrm{univ}}\colon G_\Q\to \GL_2(R)
\]
使得任意允许的提升都唯一地由 $R\to A$ 拉回得到。
\end{theorem}

\begin{proof}[证明思路]
Schlessinger 判据把可表性化为几个关于小扩张的粘合条件。
对二维 Galois 表示，
这些条件在绝对不可约的 residual 假设下可以通过群上同调验证。
本质上，
\[
  H^1(G_\Q,\mathrm{ad}^0\rhobar)
\]
控制无穷小提升，
而
\[
  H^2(G_\Q,\mathrm{ad}^0\rhobar)
\]
控制阻碍。
当局部条件选取得当时，
这些数据恰好足以构造一个泛对象 $R$。
\end{proof}

从切空间的角度看，
形变环并不是一个神秘盒子。
若 $\mathfrak m_R$ 是其极大理想，
则
\[
  \Hom_k\bigl(\mathfrak m_R/(\mathfrak m_R^2,\ell),k\bigr)
\]
可以识别为某个带局部条件的 Selmer 群。
这意味着：
$R$ 有多少个“独立方向”，
其实是一个 Galois 上同调问题。

模形式一边也有对应对象。
在与 $\rhobar$ 对应的极大理想 $\mathfrak m$ 处局部化后，
Hecke 环 $T$ 作用在模形式空间或其同调上。
由 universal modular deformation 的构造，
存在自然满射
\[
  R\twoheadrightarrow T.
\]
之所以是满射，
是因为 Hecke 特征值已经决定了模形式侧的 Galois 表示，
所以每个模提升都给出 $R$ 的一个点。
问题只剩下：
$R$ 会不会比 $T$ 更大？

\begin{theorem}[Wiles 的 $R=T$ 定理，简化表述]
\label{thm:ch10-rt}
设 $\rhobar$ 是奇、绝对不可约的 residual 表示，
并固定最小局部形变条件。
则在 Wiles--Taylor--Wiles 的情形下，
自然满射
\[
  R\twoheadrightarrow T
\]
实际上是同构。
因此每个满足这些局部条件的 $\ell$-进提升都来自模形式。
\end{theorem}

\begin{proof}[证明框架]
证明分成三层。

\textbf{第一层：数值判据。}
Wiles 的数值判据说明，
若能把 $R$ 的切空间大小与 $T$ 的 congruence ideal 大小配平，
则满射 $R\twoheadrightarrow T$ 已经足够强，可以推出同构。
其本质是一个交换代数命题：
对某类 Cohen--Macaulay 环与 complete intersection 环，
关系数与生成元数一旦恰好匹配，
满射就只能是同构。

\textbf{第二层：Selmer 群估计。}
切空间由 Selmer 群
\[
  H^1_{\mathcal L}(G_\Q,\mathrm{ad}^0\rhobar)
\]
控制。
Poitou--Tate 对偶把其维数与对偶 Selmer 群联系起来，
从而把 $R$ 的生成元数压缩成一个可计算的同调不变量。

\textbf{第三层：Taylor--Wiles 打补丁。}
真正困难处在于最小情形里，
直接比较 $R$ 与 $T$ 往往缺少自由度。
Taylor 与 Wiles 于是选择一族辅助素数 $Q_n$，满足
\[
  q\equiv 1\pmod{\ell^n},
\]
使得在这些素数添加级别结构后，
新的形变问题拥有额外变量，
而模形式侧也同步多出几乎同样数量的自由方向。
在极限上得到的大环 $R_\infty$ 与 $T_\infty$ 变成更接近幂级数环的对象，
从而可以先证明
\[
  R_\infty\cong T_\infty.
\]
最后再通过 patching 把极限层的同构下降回原始最小层，
得到 $R\cong T$。
\end{proof}

\textbf{注。}
文献常把这一节写成大量符号：
$R_Q$、$T_Q$、$\Delta_Q$、patched module、complete intersection、numerical criterion。
第一次读时最重要的不是记号，
而是记住整体图像：
\emph{辅助素数的作用不是改变问题的本质，
而是暂时把问题放到一个更有自由度的空间里，
让两边更容易比较；
比较完成后，再把多余自由度剥掉。}

把这件事想象成“打补丁”是非常贴切的：
我们并不是一次看清整个形变空间，
而是先在许多带额外级别的近似层上工作，
再把它们拼成一个大而光滑的对象，
最后回到原来的最小层。
正是这一步，
让 $R=T$ 从一句勇敢猜测变成了可操作定理。
